%% Section 4: Research Design

\section{Research Design}
\label{sec:design}

\subsection{Identification Strategy}
\label{sub:identification}

The identification relies on the intuition that the naira crunch hit all
Nigerian states simultaneously but with differential severity proportional to
each state's pre-existing dependence on physical cash. States with higher
pre-shock digital payment adoption had, in effect, a private hedge, such that when ATMs
failed, digital payment rails provided a substitute medium of exchange and
store of value. Consistent with the monetary mechanism in \citet{alvarez2022monetary}, states 
with limited digital-payment infrastructure were more exposed to the cash shortage because firms relied more heavily on 
physical currency to complete transactions, collect revenues, and pay workers.

The key identifying assumption is that, absent the crunch, high- and
low-fintech states would have followed parallel trends in enterprise outcomes.
This assumption cannot be tested directly in the firm panel, since the \NLPS{}
enterprise module was fielded only in Rounds~5, 7, and~11. It is validated
indirectly using the nighttime-light (\NTL{}) panel, which covers all 774
Nigerian LGAs for 45 consecutive pre-shock months (January~2019 to
September~2022). Because the treatment variable in both the \NTL{} pre-trend test and the
firm-panel specification is identical, the state-level fintech index
$\textit{Fintech}_s$ from \NLPS{} Round~6, the flatness of the pre-period
\NTL{} coefficients directly tests whether $\textit{Fintech}_s$ predicts
differential state-level trajectories \emph{prior} to the shock.

A secondary concern is that $\textit{Fintech}_s$ may proxy for underlying
state-level development trajectories, where richer, faster-growing states
adopted fintech earlier and also had higher economic growth irrespective of the
crunch. This concern is addressed in two ways. First, the 45-month \NTL{}
pre-period directly tests for differential trends, a formal regression of the
pre-period coefficients on calendar time yields a slope of $-0.00010$ per
month ($t = -0.42$), statistically and economically indistinguishable from
zero, confirming that high- and low-fintech states followed parallel
trajectories before the crunch. Second, enterprise fixed effects absorb all
time-invariant confounders, so the identifying variation comes from
within-enterprise changes across survey rounds rather than cross-sectional
differences.

\subsection{Primary Specification}
\label{sub:firm_spec}

For the \NLPS{} enterprise panel, the following specification is estimated:

\begin{equation}
  y_{irt} = \alpha_i + \gamma_t
    + \sum_{j \in \{R7, R11\}} \delta_j
      \bigl(\textit{Fintech}_{s(i)} \times \mathbf{1}[\text{round}_t = j]\bigr)
    + X_{irt}'\phi + \varepsilon_{irt},
  \label{eq:firm}
\end{equation}

\noindent where $y_{irt}$ is log weekly sales (or an activity or employment
indicator) for enterprise $i$ in household $r$ at round $t$; $\alpha_i$ is an
enterprise fixed effect; $\gamma_t$ absorbs round effects; and $X_{irt}$
controls for household demographics and enterprise sector. Round~5
(August~2022) is the omitted baseline. $\delta_{R7}$ captures the peak-crunch
differential and $\delta_{R11}$ captures medium-run persistence. Standard
errors are clustered at the state level.

Because the naira crunch begins for all thirty-six states and the Federal
Capital Territory simultaneously on October~26, 2022, treatment timing is
not staggered.Because treatment intensity is time invariant and enters the specification through 
its interaction with the post-crunch indicator, the estimator reduces to a standard 
two-period difference-in-differences with a continuous treatment. Under parallel trends 
and a linear dose-response relationship, the TWFE coefficient identifies the average marginal effect 
of fintech intensity on enterprise outcomes \citep{wooldridge2021two}.

That linearity assumption is probed by estimating
quintile-bin interactions and results confirm approximate linearity across the
fintech distribution (Appendix~Table~\ref{atab:alt_treatment}).

\subsection{Pre-Trend Validation}
\label{sub:pretrendvalid}

For each of the 45 pre-shock months, we estimate the coefficient on
$\textit{Fintech}_s \times \mathbf{1}[t = k]$ in a regression of
$\ln(\NTL_{lt}+1)$ on state and month-year fixed effects, where $k$ is the
month relative to February~2023 and $k=-1$ (January~2023) is the omitted
category. Under the null of parallel trends these pre-period coefficients
should be indistinguishable from zero.

The pre-period coefficients cluster around zero throughout the 45-month window.
A formal regression of the 31 estimable pre-period coefficients on relative
time yields a slope of $-0.00010$ per month ($t = -0.42$, $p = 0.68$),
confirming that high- and low-fintech states followed statistically identical
\NTL{} trajectories before the crunch. Appendix~Figure~\ref{afig:eventstudy_trend}
shows that allowing for state-specific linear trends does not materially alter
the pre-period pattern.

Appendix~Table~\ref{atab:pretrend_sensitivity} quantifies robustness to
pre-trend violations following \citet{rambachandran2023pretrends}. Setting
$M$ equal to the maximum plausible per-period gradient, the observed
pre-period slope of 0.00010 log points per month, the honest confidence
interval for the peak-crunch coefficient remains positive. The medium-run
enterprise survival result requires a pre-trend violation more than 20 times
this magnitude to be overturned, providing strong assurance that the
firm-level estimates are not artefacts of differential pre-shock trends.

\subsection{On Geographic Instruments}
\label{sub:iv}

Two geographic instruments for pre-shock fintech adoption were constructed:
road distance to the nearest financial hub (over the 2022 OpenStreetMap road
network) and terrain ruggedness (mean \textsc{tri} from the \textsc{srtm}
90-metre digital elevation model, following \citealt{hjort2019arrival}). Full
construction details and first-stage results are in
Appendix~Table~\ref{atab:iv_firststage}. Both instruments fail the relevance
condition ($F = 2.14$, $p = 0.13$). The road-distance coefficient has the
wrong sign such thatstates farther from financial hubs score \emph{higher} on
the fintech index, suggesting a leapfrog dynamic where the most financially
excluded states adopted mobile money most rapidly between 2014 and
2022 \citep{suri2016mobile}; Appendix~Table~\ref{atab:leapfrog} documents
this pattern. Identification therefore rests entirely on
equation~\eqref{eq:firm} with the 45-month pre-shock \NTL{} window as the
parallel-trends validation.
